\section{Related Literature}\label{sec:literature}

This paper connects three literatures: the real effects of cartel enforcement, the economics of worker mobility and knowledge transfer, and the econometrics of staggered difference-in-differences.

\subsection{Real effects of cartel enforcement}

A growing body of work traces the consequences of antitrust enforcement beyond the direct penalties imposed on convicted firms. \citet{szerman2023employee} exploits a debarment-policy change in Brazilian procurement to show that blacklisted firms lose nearly half their employees---a substantial labor-market disruption. Our paper asks the next question: where do those employees go, and what happens when they arrive?

On the product-market side, the enforcement literature has documented that cartel breakdowns reduce prices \citep{connor2005price} and increase entry \citep{harrington2008detecting_effects}. But these studies treat the transition from collusion to competition as a within-market event, without tracking the inter-firm reallocation of human capital that accompanies it. We show that this reallocation is a channel through which enforcement reshapes the competitive landscape: the firms that absorb ex-cartel workers become stronger competitors in the same markets.

\subsection{Worker mobility and knowledge transfer}

Labor economists have established that workers are conduits for inter-firm knowledge transfer. \citet{stoyanov2012worker} showed that productivity spillovers flow between Danish firms through mobile workers, with larger effects when the productivity gap between sender and receiver is wide. \citet{serafinelli2019good} documented similar patterns in Italy: workers arriving from high-productivity firms raise wages and output at their new employers. \citet{jarosch2021learning} developed a model in which worker reallocation shapes firm-level human capital through on-the-job learning.

These papers establish that worker flows transmit \emph{productive} knowledge. But can they also transmit \emph{collusive} knowledge? \citet{lamichhane2024collusion} provide the first direct evidence, showing in a laboratory experiment that worker mobility between Bertrand competitors facilitates collusive pricing: workers carry pricing-strategy information, and firms that exchange employees coordinate on higher prices. \citet{herrera2025collusion} extend the logic to shared leadership: common directors and executives raise the probability of product-market collusion by 11~percentage points. Strikingly, inter-firm worker flows \emph{fall} by 20\% once board interlocks are established, suggesting that leadership ties and worker flows are substitute channels for coordination.

Our paper provides the first field test of whether the Lamichhane--McGee mechanism operates in practice after enforcement-driven labor reallocation. The answer is no: the workers who move are predominantly operational staff, the receiving firms' win rates rise but their prices do not, and the effect is concentrated in capacity deepening rather than coordination.

\subsection{Procurement and thin labor markets}

Two papers have linked RAIS to Brazilian procurement, establishing the data infrastructure we build on. \citet{barbosa2020revolving} track revolving-door movements between public agencies and private suppliers, finding that supplier-to-administration transitions are detrimental to procurement efficiency. \citet{szerman2023employee} uses the RAIS--procurement linkage to estimate worker-level costs of debarment. Neither paper examines the competitive effects on receiving firms.

The thin-market structure of S\~{a}o Paulo procurement is central to our identification. In markets with few specialized suppliers, a cartel firm's contraction mechanically concentrates the displaced workers in a small set of competitors. This concentration produces a ``dose'' of operational capacity that is large enough to shift win rates---an effect that would be diluted in thick markets with many potential absorbers.

\subsection{Staggered DiD and TWFE bias}

Our setting features staggered treatment: different firms absorb ex-cartel workers in different years as cartel conduct periods end at different times. \citet{goodman2021difference} showed that the standard TWFE estimator can be biased in staggered settings because already-treated units serve as implicit controls for later-treated units. \citet{cengiz2019effect} proposed the stacked-cohort approach we adopt, and \citet{callaway2021difference} and \citet{borusyak2024revisiting} developed alternative estimators with similar properties.

Our paper provides a vivid illustration of TWFE bias: the naive TWFE specification produces a highly significant price ``contagion'' effect ($+0.26$ log points, $t > 5$) that vanishes entirely under the stacked estimator ($+0.05$, insignificant). Had we relied on TWFE alone, the paper would have reported a false positive with direct policy implications---a cautionary example for the applied literature.
